\begin{tabular}{l r r r r r}
\toprule
\multicolumn{6}{l}{\textit{Panel A: Coverage}} \\
\midrule
\multicolumn{5}{l}{Firms with debt data} & 409 / 451 (90.7\%) \\
\multicolumn{5}{l}{Events covered} & 120 / 159 (75.5\%) \\
\multicolumn{5}{l}{Firm-event observations (h=0)} & 37,114 / 64,851 (57.2\%) \\
\midrule
\multicolumn{6}{l}{\textit{Panel B: Debt-to-assets distribution at h=0 (winsorized 1/99 pct)}} \\
\midrule
N & Mean & SD & p25 & Median & p75 \\
\midrule
37,114 & 0.244 & 0.175 & 0.105 & 0.233 & 0.350 \\
\bottomrule
\multicolumn{6}{p{\linewidth}}{\footnotesize \textit{Notes:} Matched SEC XBRL subsample. sample\_debt = sample\_preferred AND has\_debt\_data. The debt-to-assets ratio (debt\_assets\_lag\_w) uses lagged total debt over lagged total assets from each firm's most recent 10-Q prior to the FOMC event, winsorized at the 1st/99th percentiles. The high\_debt event-time dummy splits firms by the within-event median (not shown). Panel A denominators are the preferred sample (firms, events, and firm-event h=0 observations within sample\_preferred==1). Sample: Stata baseline window 2005-2024 (reproduces current thesis tables). Built by 00\_master\_data\_pipeline\_and\_data\_section.ipynb (Python; full available panel; Stata applies \$LP\_DATE\_FILTER separately).} \\
\end{tabular}
